% -------------------------------------------------
% Case histories
% -------------------------------------------------

\chapter{Case Histories}
\label{chap:case-histories}

This chapter contains fictional clinical scenarios for practising
patient-centred history taking. The actor should read the complete
brief before the interview but should reveal information only when it
is asked for naturally.

\section{Case 1: Oesophageal Cancer}
\label{sec:case-oesophageal-cancer}

\begin{tabularx}{\textwidth}{@{}>{\bfseries}p{0.31\textwidth}Y@{}}
Name
    & David Morgan \\
Age
    & 74 years \\
Occupation
    & Retired bus driver \\
Presenting concern
    & Increasing difficulty swallowing \\
Setting
    & General practice appointment \\
Intended condition
    & Oesophageal cancer \\
\end{tabularx}


\subsection*{Patient identity}

You are David Morgan, a 64-year-old retired bus driver. You live with
your partner, Helen, and are normally independent. You have come to see the doctor because swallowing has become progressively more difficult.

You are polite and cooperative but visibly worried. You have delayed
making the appointment because you hoped the problem would settle by
itself.


\subsection*{History of the presenting complaint}

\begin{description}
    \item[Onset:]
    The problem began approximately four months ago. There was no
    particular event, illness or injury when it started.

    \item[Progression:]
    The difficulty has gradually become worse. It was initially
    occasional but now occurs during most meals.

    \item[Solid food:]
    At first, bread and meat felt as though they were sticking. You now
    avoid most solid food because it is difficult to swallow.

    \item[Soft food:]
    During the last few weeks, softer foods such as mashed potato have
    sometimes also been difficult to swallow.

    \item[Liquids:]
    You can still drink liquids. Occasionally they seem to pass slowly,
    but they do not make you cough or choke.

    \item[Location:]
    Food feels as though it sticks behind the lower part of your
    breastbone. You have no difficulty starting a swallow.

    \item[Pain:]
    Swallowing is uncomfortable when food sticks, but it is not usually
    painful.

    \item[Regurgitation:]
    On three or four occasions, a small amount of undigested food has
    come back into your mouth after it became stuck. You have not had
    forceful vomiting.

    \item[Weight loss:]
    You have unintentionally lost approximately \SI{6}{\kilogram} in
    four months. Your weight has fallen from about \SI{82}{\kilogram}
    to \SI{76}{\kilogram}.

    \item[Appetite:]
    Your appetite is still present. You eat less because meals take a
    long time and you are worried about food becoming stuck.

    \item[Fatigue:]
    You have felt more tired than usual during the last month, although
    you can still manage your usual daily activities.

    \item[Heartburn:]
    You have experienced intermittent heartburn for several years,
    usually after a large or spicy meal. You take an over-the-counter
    antacid when necessary. The heartburn itself has not recently
    become worse.
\end{description}

\subsection*{Associated symptoms}

Reveal the following information only if the student asks:

\begin{itemize}
    \item No coughing or choking when beginning to swallow.
    \item No pain in the mouth or throat.
    \item No persistent nausea or vomiting.
    \item No blood in vomit.
    \item No black, tar-like stools and no visible blood in the stool.
    \item No severe chest pain.
    \item No persistent abdominal pain.
    \item No change in voice.
    \item No persistent cough or breathlessness.
    \item No weakness, facial drooping or other neurological symptoms.
    \item No fevers or night sweats.
    \item No noticeable lumps in the neck.
\end{itemize}

\subsection*{Past medical and surgical history}

\begin{description}
    \item[Medical history:]
    You have experienced intermittent acid reflux for several years.
    You were diagnosed with high blood pressure three years ago.

    \item[Previous investigation:]
    You have never had a camera test of your oesophagus or stomach.

    \item[Surgical history:]
    You had your appendix removed in your twenties. You have had no
    other operations.

    \item[Previous cancer:]
    You have never previously been diagnosed with cancer.
\end{description}

\subsection*{Medication and allergies}

\begin{description}
    \item[Regular medication:]
    Ramipril \SI{5}{\milli\gram} once daily for high blood pressure.

    \item[Other medication:]
    An over-the-counter antacid once or twice a week when required. Cocaine user.

    \item[Allergies:]
    No known drug allergies.
\end{description}

\subsection*{Family history}

Your father died following a heart attack at the age of 71. Your mother
died at the age of 83 following a stroke.

You have one older sister, who is generally well. There is no known
family history of oesophageal or stomach cancer.

\subsection*{Social history}

\begin{description}
    \item[Home:]
    You live with your partner, Helen, in a two-storey house. You are
    independent with personal care, cooking and household activities.

    \item[Occupation:]
    You retired from bus driving four years ago.

    \item[Smoking:]
    You smoked approximately 15 cigarettes per day from the age of 20
    until the age of 59. You stopped smoking five years ago.

    \item[Alcohol:]
    You usually drink three or four pints of beer over the course of a
    weekend, amounting to approximately 6--8 units each week.

    \item[Diet:]
    You previously ate a varied diet. You now mainly have soup, yoghurt,
    porridge and other soft foods. You drink water with meals to help
    food pass.

    \item[Recreational drugs:]
    You have never used recreational drugs.
\end{description}

\subsection*{Effect on daily life}

Meals now take much longer than usual, and you often leave food
unfinished. You have stopped eating meat and bread and have begun
avoiding meals in restaurants because you are embarrassed that food
might become stuck.

Helen has noticed the weight loss and encouraged you to make this
appointment. You remain independent but have less energy for gardening
and walking.

You are worried that you will continue losing weight or eventually
become unable to swallow anything.

\subsection*{Ideas, concerns and expectations}


\begin{description}
    \item[Ideas:]
    You initially thought the problem was caused by indigestion. You now
    wonder whether something is narrowing your \enquote{food pipe}.

    \item[Concerns:]
    Your main fear is cancer, particularly because the problem is
    getting worse and you have lost weight. You are also worried about
    how Helen would cope if you became seriously ill.

    \item[Expectations:]
    You want to know what could be causing the problem and whether you
    will need a scan or a camera test. You hope that the cause is
    treatable.
\end{description}


\subsection*{Questions the actor may ask}

If there is time, ask one or two of the following:

\begin{itemize}
    \item \enquote{What could be causing this?}
    \item \enquote{Do you think it might be cancer?}
    \item \enquote{What tests will I need?}
    \item \enquote{Is there anything I can do about the weight loss?}
    \item \enquote{How quickly will I be seen?}
\end{itemize}

The actor should accept an explanation that further assessment is
needed and that the student cannot confirm a diagnosis from the history
alone.

\clearpage

\subsection*{Tutor information}

This is a fictional teaching case. The pattern of progressive
swallowing difficulty, beginning with solid food and later affecting
softer food, suggests a mechanical obstruction of the oesophagus.
Unintentional weight loss and increasing symptoms are concerning
features that require prompt investigation.

The intended diagnosis is oesophageal cancer. It should not be treated
as confirmed during the interview because no investigation results
have been provided.

\subsection*{Features the student should identify}

\begin{itemize}
    \item Progressive difficulty swallowing over four months.
    \item Initial difficulty with solid food, followed by softer food.
    \item Sensation of food sticking behind the breastbone.
    \item Unintentional loss of \SI{6}{\kilogram}.
    \item Reduced food intake despite a preserved appetite.
    \item Increasing fatigue.
    \item Longstanding reflux symptoms.
    \item Previous smoking history.
    \item Significant anxiety about cancer and the effect of illness on
    his partner.
\end{itemize}

\subsection*{Expected clinical approach}

The student should recognise that progressive dysphagia with
unintentional weight loss requires urgent clinical assessment. A good
consultation should include:

\begin{itemize}
    \item clarification of whether the difficulty is with initiating a
    swallow or with food passing through the oesophagus;
    \item assessment of progression, regurgitation, pain, bleeding and
    nutritional consequences;
    \item relevant gastrointestinal, respiratory and neurological
    symptoms;
    \item medical, medication, family and social histories;
    \item assessment of smoking and alcohol exposure;
    \item exploration of ideas, concerns and expectations;
    \item acknowledgement of the patient's concern without confirming
    or excluding cancer prematurely;
    \item explanation that prompt investigation is required, commonly
    including upper gastrointestinal endoscopy and biopsy where
    indicated; and
    \item consideration of nutritional support and safety-netting while
    investigation is arranged.
\end{itemize}

\subsection*{Possible discussion points}

\begin{enumerate}
    \item Which features make this an urgent presentation?
    \item What alternative diagnoses could produce progressive
    dysphagia?
    \item What physical examination would be appropriate?
    \item What initial investigations and referrals should be
    considered?
    \item How should the possibility of cancer be discussed without
    implying that the diagnosis is already confirmed?
    \item What nutritional and psychosocial support may be helpful?
\end{enumerate}

\subsection*{Teaching note}

The precise referral pathway, investigations and management should
follow current local and national guidance. No tumour location,
histological type or stage is specified in this scenario, and these
details should not be assumed.
